\input{algebra1.tex}
\title[Algebra 1.]{Polinomok faktorizációja}
\subtitle{}
%\subtitle{Mátrixok bevezetése}
\date[4. hét]{2026.~október 5-én}
\begin{document}
\openingframes{Tartalom}
\section{Polinom faktorizáció}
\frame{
Kicsi korunk óta sulykolják belénk, hogy minden egész szám előáll, méghozzá lényegében csak egyféleképpen
prímek szorzataként.
Ha ismerjük két szám prímtényezős előállítását, akkor nagyon könnyű megmondani a két szám legkisebb közös többszörösét,
vagy a legnagyobb közös osztóját.
Evvel a probléma csak annyi,
hogy nagyon nehéz megmondani két, esetleg jó nagy, szám prímtényezős előállítását,
így még az egész számok gyűrűjében is az Euklideszi-algoritmus a megfelelő, véges sok lépésben,
végrehajtható módszer a legnagyobb közös osztó és a legkisebb közös többszörös konkrét felírására.

Ebben a szakaszban azt mutatjuk meg, hogy prímtényezős előállításról szóló tétel a polinomok gyűrűjében is igaz marad.
Természetesen konkrét algoritmust nem adunk, hiszen ilyen még számokra sem igen van.%
\begin{definition}[reducibilis polinom]\index{reducibilis polinom}\index{irreducibilis polinom}
	Egy polinomot \emph{reducibilisnek} mondunk,
	ha előáll mint két legalább elsőfokú polinom szorzata.
	Egy nem reducibilis polinomot \emph{irreducibilisnek} nevezünk.
\end{definition}
}
\frame{
Világos, hogy minden legfeljebb elsőfokú polinom tetszőleges test felett irreducibilis.
A magasabb fokú polinomok esetében a probléma nagyban függ a testtől is,
ahonnan a polinom együtthatói származnak.
A következő állítás viszont minden test mellett igaz.
\begin{proposition}[polinom faktorizáció]
	Tetszőleges test feletti polinomgyűrűben,
	minden (normált) polinom előáll mint (normált) irreducibilis polinomok szorzata.
    \label{pr:polfact}
\end{proposition}
%\begin{proof}
%	Az előállítandó polinom foka szerinti teljes indukció.
%	Elsőfokú polinom maga irreducibilis.
%
%	Most tegyük fel, hogy az állítás igaz $n$-nél alacsonyabb fokú polinomokra
%	és lássuk be $\deg p=n$ mellett, ahol $n>1$.
%	Ha $p$ irreducibilis, akkor megint készen vagyunk.
%	Ha $p=f g$ valamely $\deg f\geq 1$ és $\deg g\geq 1$ mellett,
%	akkor $\deg f<n$ és $\deg g<n$.
%	Az indukciós feltevés szerint $f$ és $g$, emiatt $p=fg$ is előáll irreducibilis polinomok szorzataként.
%\end{proof}
}
\section{Prim tulajdonság}
\frame{
Érdemes látni,
hogy ha a $d$ polinom egy irreducibilis $p$ polinom osztója,
akkor csak
\[
    \deg d=0, \text{ vagy } \deg d=\deg p
\]
lehetséges.

Ezért \emph{egy tetszőleges $f$ polinomra igaz, hogy a $p$ irreducibilis polinomra
vagy $p|f$ vagy $p$ és $f$ relatív prímek.}

Ugyanis,
ha az $f$ és $p$ polinomok $d$ legnagyobb közös osztója nem $1$,
akkor csak $c\cdot d=p$ lehetséges valamely $c\in\mathbb{F}$ konstans mellett, 
ergo $p|f$.
}
\frame{
\begin{proposition}[prím tulajdonság]\index{prím polinom}
	Legyen $p\in\mathbb{F}\left[ t \right]$ irreducibilis polinom, amelyre
	$p|(f_1\cdot\ldots\cdot f_n)$, valamely $f_j\in\mathbb{F}\left[ t \right]$ polinomokra, ahol $j=1,\ldots,n\geq 1$.
	Ekkor létezik $1\leq j\leq n$, amelyre $p|f_j$.\qedhere
\end{proposition}
%\begin{proof}
%	A polinomok $n$ száma szerinti indukció.
%	Az $n=1$ eset semmitmondó módon teljesül.
%	Tegyük fel, hogy igaz az állítás $n$-nél kevesebb polinomra,
%	és lássuk be $n$-re. Itt $n\geq 2$.
%	Induljunk ki tehát abból, hogy
%	\[
%		p|\left( f_1\cdots f_{n-1} \right)\cdot f_n
%	\]
%	Ha $p$ osztója lenne az $f_1\cdots f_{n-1}$ szorzatnak,
%	akkor az indukciós feltevés szerint készen is lennénk.
%	Ha $p$ nem osztója a szorzatnak,
%	akkor $p$ irreducibilis volta miatt relatív prímek.
%	Ekkor a \ref{pr:rprim}. állítás szerint $p|f_n$.
%\end{proof}
Az is igaz, hogy ha egy polinom reducibilis, akkor az nem prím.

Összeségében látjuk tehát,
hogy 
\emph{a polinomok irreducibilitása és prím tulajdonsága azonos fogalmak.}
}
\section{A faktorizáció egyértelműsége}
\frame{
Az irreducibilis polinomok prím tulajdonsága segítségével a polinomok faktorizáció egyértelműségét is igazolhatjuk.
\begin{proposition}
	Legyen $p \in\mathbb{F}\left[ t \right]$ egy legalább elsőfokú normált polinom.
	Ekkor ezek sorrendjétől eltekintve egyértelműen léteznek legalább elsőfokú, normált, irreduciblis $q_1,\dots,q_s\in\mathbb{F}\left[ t \right]$ polinomok,
	hogy $p=q_1\cdot\ldots\cdot q_s$.
\end{proposition}
A ,,sorrendtől eltekintés'' alatt azt értjük, hogy ha $p$ előáll
\[
	p_1\cdot\ldots\cdot p_s=p=q_1\cdot\ldots\cdot q_r
\]
legalább elsőfokú, normált, irreducibilis polinomok szorzataként,
akkor $s=r$ és a $p_1,\ldots,p_s$ polinomok alkalmas átindexelése után $p_j=q_j$,
minden $j=1,\ldots,s$ mellett.
%\begin{proof}
%	\Aref{pr:polfact}. állítás
%	szerint $p$ előáll, mint normált, irreducibilis polinomok szorzata.
%	Itt legalább egy polinom legalább elsőfokú, hiszen $p$ legalább elsőfokú.
%	Ha van a szorzat tagjai közt nulladfokú, akkor az csak 1 lehet a normáltság szerint,
%	ezért egyszerűen elhagyható.
%	Meggondoltuk tehát, hogy $p$ legalább elsőfokú, normált, irreducibilis polinomok szorzata.
%
%	Az unicitást,
%	a felbontandó polinom fokszáma szerinti indukcióval igazoljuk:
%	Ha a polinom elsőfokú, akkor a felbontás egyértelműsége is nyilvánvaló.
%	Tegyük fel, hogy igaz az egyértelműség $n$-nél alacsonyabb fokszámú polinomokra,
%	és tegyük fel, hogy $\deg p=n>1$.
%	Nézzünk két lehetséges előállítást
%	\[
%		p_1\cdots p_s=p=q_1\cdots q_r,
%	\]
%	ahol $p_1,\dots,p_s,q_1,\dots,q_r$ legalább elsőfokú, normált, irreducibilis polinomok.
%	A $q_1$ polinom osztója a jobboldalnak, ezért a baloldalnak is.
%	A prím tulajdonság miatt, \ref{pr:rprim}. állítás, $q_1$ osztója az egyik baloldali polinomnak.
%	Alkalmas átindexelés után feltehető, hogy $q_1|p_1$.
%	No de, $p_1$ is irreducibilis, és $\deg q_1\geq 1$ miatt csak $\deg q_1=\deg p_1$ lehetséges,
%	tehát a normáltság szerint $q_1=p_1$.
%	A polinomgyűrű nullosztó mentessége szerint az első polinomokkal egyszerűsíthetünk, ergo
%	\[
%		p_2\cdots p_s=q_2\cdots q_r,
%	\]
%	is fennáll.
%	A fenti polinom már $n$-nél alacsonyabb fokú,
%	így az indukciós feltevés szerint $s-1=r-1$,
%	és alkalmas átindexelés után minden $j=2,\dots,s$ esetén is teljesül a $p_j=q_j$ egyenlőség.
%\end{proof}
}
\frame{
Az egész szakaszt összefoglalhatjuk így is:
\begin{proposition}
	Minden legalább elsőfokú normált polinomhoz léteznek
	--
	méghozzá sorrendjüktől eltekintve egyértelműen léteznek
	--
	olyan $q_1,\dots,q_s$ egymástól különböző, normált, irreducibilis polinomok;
	és olyan $n_1,\dots,n_s$ pozitív egészek;
	amelyekre
	\[
		p\left( t \right)
		=
		q_1^{n_1}\left( t \right)\cdots q_s^{n_s}\left( t \right).\qedhere
	\]
\end{proposition}
}
\end{document}
